% Generated draft: requires mathematical and source review.
\documentclass[11pt]{article}
\usepackage[margin=25mm]{geometry}
\usepackage{fontspec}
\setmainfont{Latin Modern Roman}
\usepackage{amsmath,amssymb,mathtools,unicode-math}
\setmathfont{Latin Modern Math}
\usepackage{enumitem,needspace,xurl,hyperref}
\setlength{\parindent}{0pt}
\setlength{\parskip}{5pt}
\setlength{\emergencystretch}{3em}
\newcommand{\N}{\mathbb{N}}
\newcommand{\Z}{\mathbb{Z}}
\newcommand{\Q}{\mathbb{Q}}
\newcommand{\R}{\mathbb{R}}
\newcommand{\C}{\mathbb{C}}
\newcommand{\F}{\mathbb{F}}
\newcommand{\A}{\mathbb{A}}
\newcommand{\PP}{\mathbb{P}}
\newcommand{\E}{\mathbb{E}}
\newcommand{\eps}{\varepsilon}
\newcommand{\dd}{\,\mathrm{d}}
\begin{document}
\section*{UnsolvedMath ID: 30005171}
\textbf{Source problem code:} catalogue entry
\subsection{Definitions and mathematical statement}
Let \(p\) be a prime integer. The question concerns smooth projective curves that arise as smooth limits of smooth plane curves of degree \(p\). The supplied record does not specify the ground field, the precise class of families, or whether finite base change is allowed. Consequently, the following is a provisional unpacking of the question, with those conventions still requiring confirmation.

Under the usual one-parameter interpretation, a smooth limit is a smooth projective curve \(C\) occurring as the special fiber \(\mathcal C_0\) of a flat proper family of curves \(\mathcal C\to T\) over a pointed integral curve \((T,0)\), whose geometric general fiber is isomorphic to a smooth degree-\(p\) curve in \(\PP^2\). The ground field and any necessary restrictions on \(T\) remain unspecified.

An effective Cartier divisor on a scheme \(Y\) is a closed subscheme whose ideal sheaf is locally generated by one non-zero-divisor. A degeneration of \(\PP^2\) is provisionally understood as a flat projective family of surfaces \(\pi\colon\mathcal X\to S\) over a pointed one-parameter base \((S,0)\), with geometric general fiber isomorphic to \(\PP^2\). The basic relative \(\Q\)-Gorenstein condition is that the relative canonical divisor \(K_{\mathcal X/S}\) is \(\Q\)-Cartier, meaning that some positive integral multiple is Cartier. The source does not specify whether a stronger condition involving reflexive powers of the relative dualizing sheaf and their compatibility with base change is intended. For a normal \(\Q\)-Gorenstein variety \(Y\), log terminality, with zero boundary, means that every exceptional prime divisor \(E\) over \(Y\) has discrepancy \(a(E,Y)>-1\), in a setting where these discrepancies are defined.

\textbf{Question, with unresolved conventions.} Is it true that, for every prime \(p\) and every smooth limit \(C\) of smooth plane curves of degree \(p\), there exists a log-terminal \(\Q\)-Gorenstein degeneration of \(\PP^2\) in which \(C\) occurs as a Cartier divisor?

One natural precise interpretation of the desired conclusion is the existence of a family \(\pi\colon\mathcal X\to S\) as above and a closed immersion
\[
 i\colon C\hookrightarrow X_0,\qquad X_0=\pi^{-1}(0),
\]
such that \(i(C)\) is an effective Cartier divisor on the normal log-terminal surface \(X_0\). This interpretation requires source confirmation: the supplied wording does not determine whether log terminality is required of the special surface, the total space, or both. It also does not determine whether only an embedding of the abstract curve \(C\) is required, or whether the given family of plane curves must extend, possibly after base change, to a family of Cartier divisors on the surface degeneration. No conclusion about these stronger requirements is asserted here. The prime-degree question includes \(p=2\) and \(p=3\); the restriction \(p>2\) below belongs only to the supplied background result.

\textbf{Supplied background, not independently verified.} A Markov number is a positive integer appearing in a positive integral triple \((a,b,c)\) satisfying
\[
 a^2+b^2+c^2=3abc.
\]
The excerpt states as ``Theorem A'' that, for every Markov number \(d>2\), there is a family of smooth plane curves of degree \(d\) with a smooth projective nonplanar limit; in particular, this applies to every prime Markov number \(p>2\). Here nonplanar is understood to mean that the limiting abstract curve admits no embedding as a smooth plane curve. The excerpt also attributes to Griffin an example of smooth plane quintics with a hyperelliptic, hence nonplanar, limit. These background assertions concern failure of planarity and do not themselves settle the Cartier-divisor question. The supplied excerpt is truncated and contains an incomplete reference to another conjecture, so its missing hypotheses cannot be recovered from this record.
\subsection{Short English statement}
For every prime degree, must every smooth curve obtained as a limit of smooth plane curves occur as a Cartier divisor in a log-terminal \(\Q\)-Gorenstein degeneration of \(\PP^2\)? The field, degeneration conventions, and required compatibility with the original family need clarification. The supplied record labels the problem partially solved; that assessment has not been independently verified.
\subsection{Sources}
\begin{itemize}
\item Supplied local UnsolvedMath registry record, permanent ID \(30005171\), source problem code \texttt{catalogue entry}, title \emph{Stable Limits of Smooth Plane Curves}; registry locations identified by the project are \texttt{data/unsolvedmath/problems.json} and \texttt{catalog/problems.json}. This draft uses the supplied record and does not assert independent inspection of those files.
\item Supplied report citation: \emph{Algebraic Geometry: Moduli Spaces, Birational Geometry and Derived Aspects (2023)}, Oberwolfach Reports, pp.\ 1805--1864, DOI \texttt{10.4171/owr/2022/32}. Supplied literature-source label: ``Oberwolfach Report 32/2022.'' \url{https://doi.org/10.4171/owr/2022/32}. The supplied title year and report designation differ and require bibliographic confirmation.
\item Supplied literature reference: DeVleming--Stapleton, \emph{Smooth limits of plane curves}. \url{https://doi.org/10.5802/jep.263}. The record also attributes a quintic example to Griffin, reference [8] in its excerpt, without supplying that reference's bibliographic entry. No individual proposer is identified.
\item Supplied literature assessment dated \textbf{2026-08-22}: the problem is labeled ``partially solved''; the assessment says that nonplanar smooth limits occur in Markov prime degrees while the claimed general classification remains open. This is supplied dated curation, not an independently verified literature assessment. The record explicitly requests independent verification of time-sensitive status and cited claims.
\item The registry labels the statement \texttt{corrected\_verified} and reports that the official Oberwolfach report was checked and the question isolated from surrounding exposition. These are supplied provenance claims; no source visit or independent verification is asserted in this draft.
\end{itemize}
\end{document}
